%% =====================================================================
%%  Bienvenue : la promesse de la méthode, les niveaux, le rythme
%% =====================================================================
\liminaire{bienvenue}{Bienvenue !}{\enb{Welcome!} — Tout ce qu'il faut savoir avant de commencer}

{\fontsize{11.4}{15.6}\selectfont
Vous voulez apprendre l'anglais et vous partez de zéro ? Vous l'avez étudié il y a longtemps
et il ne vous en reste presque rien ? Cette méthode est faite pour vous. Elle ne suppose
\textbf{aucune connaissance préalable} : chaque mot, chaque règle et chaque son est expliqué
en français, simplement, puis entraîné jusqu'à devenir naturel.\par\vspace{3mm}

Vous suivrez Léa, une Française qui vient de s'installer à New York pour son travail. Avec
elle, vous apprendrez à vous présenter, à parler de votre quotidien, à vous débrouiller en
ville, à inviter des amis et à raconter vos week-ends. Pas de listes à rallonge : uniquement
l'anglais dont vous aurez réellement besoin, dans l'ordre où vous en aurez besoin.\par}

\otIntertitre{À la fin de ce tome, vous saurez\ldots}
\begin{paire}
\begin{tcolorbox}[otBase,colback=otNavyTint,colframe=otNavyTint]
\begin{list}{\otCheck}{\setlength{\leftmargin}{7mm}\setlength{\labelwidth}{4mm}\setlength{\labelsep}{3mm}
  \setlength{\itemsep}{1.6mm}\setlength{\topsep}{0pt}\setlength{\parsep}{0pt}}\fontsize{10.4}{13}\selectfont
\item vous présenter et présenter vos proches ;
\item épeler votre nom, donner votre numéro et votre adresse ;
\item parler de votre journée, de vos goûts, de votre logement ;
\item demander votre chemin, faire des achats, commander au restaurant ;
\end{list}
\end{tcolorbox}
\begin{tcolorbox}[otBase,colback=otNavyTint,colframe=otNavyTint]
\begin{list}{\otCheck}{\setlength{\leftmargin}{7mm}\setlength{\labelwidth}{4mm}\setlength{\labelsep}{3mm}
  \setlength{\itemsep}{1.6mm}\setlength{\topsep}{0pt}\setlength{\parsep}{0pt}}\fontsize{10.4}{13}\selectfont
\item dire ce que vous savez faire et ce qui se passe en ce moment ;
\item parler de la météo, des dates, de vos projets ;
\item raconter ce que vous avez fait hier ou pendant un voyage ;
\item comprendre l'essentiel d'une conversation simple et lente.
\end{list}
\end{tcolorbox}
\end{paire}

\otIntertitre{Le niveau A1, première marche d'un parcours en cinq tomes}
Les niveaux de langue sont définis par le \textbf{CECRL} (Cadre européen commun de référence
pour les langues), utilisé dans toute l'Europe par les écoles, les entreprises et les
certifications. Ce tome vous mène du niveau zéro au niveau \textbf{A1} ; les tomes suivants
vous accompagneront jusqu'au niveau C1.\par\vspace{4mm}

\noindent\begin{tikzpicture}
  \foreach \n/\niv/\nom/\t/\c in {1/A1/Débutant/Tome 1/otGreen,2/A2/Élémentaire/Tome 2/otGreenDark,
      3/B1/Intermédiaire/Tome 3/otNavy!70,4/B2/Avancé/Tome 4/otNavy,5/C1/Autonome/Tome 5/otNavyDark,6/C2/Maîtrise/ /otMuted!45}{%
    \pgfmathsetmacro{\x}{(\n-1)*28.5}
    \pgfmathsetmacro{\h}{10+(\n-1)*7}
    \fill[\c,rounded corners=1.2mm] (\x mm,0) rectangle ++(26.5mm,\h mm);
    \node[text=white,font=\otBold\fontsize{15}{16}\selectfont] at (\x mm+13.25mm,\h mm-5.6mm) {\niv};
    \node[anchor=north,font=\otSemi\fontsize{8.4}{10}\selectfont,text=otInk] at (\x mm+13.25mm,-1.6mm) {\nom};
    \node[anchor=north,font=\otRegular\fontsize{7.8}{9.4}\selectfont,text=otMuted] at (\x mm+13.25mm,-5.6mm) {\t};}
  \node[anchor=south,fill=otOrange,text=white,rounded corners=1.4mm,inner xsep=2.4mm,inner ysep=1.2mm,
    font=\otSemi\fontsize{8}{9}\selectfont] (ici) at (13.25mm,17mm) {Vous êtes ici};
  \draw[otOrange,line width=1pt,-{Stealth[length=2mm]}] (ici.south) -- (13.25mm,11mm);
  \node[anchor=north west,font=\otRegular\fontsize{7.8}{9.4}\selectfont,text=otMuted,align=left] at (0mm,45mm)
    {Le niveau C2, celui des traducteurs et\\ des interprètes, dépasse le cadre de la collection.};
\end{tikzpicture}

\clearpage
\otIntertitre{Votre rythme : une unité par semaine}
Une unité demande environ \textbf{trois heures} de travail. Le plus efficace est de les répartir
en \textbf{six séances de 30 minutes}, que vous retrouverez sur la page d'ouverture de chaque
unité. À ce rythme, vous terminerez le tome en cinq mois environ, révisions comprises.\par\vspace{4mm}

\noindent\begin{tikzpicture}
  \foreach \j/\d/\s/\c in {1/Lun./Dialogues/otNavy,2/Mar./Vocabulaire/otGreenDark,3/Mer./Grammaire/otViolet,
      4/Jeu./Sons et culture/otOrangeDark,5/Ven./Exercices/otNavy,6/Sam./Bilan/otGreenDark,7/Dim./Repos ou révision/otMuted}{%
    \pgfmathsetmacro{\x}{(\j-1)*24.4}
    \fill[\c!12,rounded corners=1.4mm] (\x mm,0) rectangle ++(22.6mm,20mm);
    \fill[\c,rounded corners=1.4mm] (\x mm,13mm) rectangle ++(22.6mm,7mm);
    \node[text=white,font=\otSemi\fontsize{8.4}{10}\selectfont] at (\x mm+11.3mm,16.5mm) {\d};
    \node[text width=20mm,align=center,font=\otMedium\fontsize{7.8}{9.4}\selectfont,text=otInk] at (\x mm+11.3mm,6.5mm) {\s};}
\end{tikzpicture}

\otIntertitre{Les six clés de la réussite}
\begin{tcbraster}[raster columns=3,raster equal height=rows,raster column skip=3.5mm,raster row skip=3.5mm,
  raster before skip=1mm,raster after skip=2mm]
\otCartePrg{horloge}{otNavy}{Un peu chaque jour}{Trente minutes par jour valent mieux que trois heures le dimanche : la régularité fait tout.}
\otCartePrg{micro}{otGreenDark}{À voix haute}{Lisez, répétez, répondez toujours à voix haute. On apprend à parler en parlant.}
\otCartePrg{casque}{otViolet}{Écouter d'abord}{Écoutez chaque dialogue avant de le lire : votre oreille s'habitue aux vrais sons.}
\otCartePrg{livre}{otOrangeDark}{Réviser au bon moment}{Revoyez les mots le lendemain, puis trois jours après, puis une semaine après.}
\otCartePrg{etoile}{otTeal}{Oser se tromper}{Les erreurs font partie de l'apprentissage. Les corrigés vous expliquent chacune d'elles.}
\otCartePrg{globe}{otRedDark}{L'anglais partout}{Chansons, séries en version originale, applications : chaque contact compte.}
\end{tcbraster}

\begin{astuce}[Avant de commencer]
Lisez les pages suivantes : le \textbf{mode d'emploi} (p.~\pageref{lim-mode}), la façon d'\textbf{écouter
les audios} (p.~\pageref{lim-audios}) et le \textbf{guide des sons} (p.~\pageref{lim-sons}). Dix minutes
bien investies : vous saurez exactement comment travailler chaque page.
\end{astuce}
